% Nouveau EGA V (cote 25): archive cote 25, site page 167 (PDF page 168).
% typescript. Draft compared: none.
% AI-transcribed by Claude Opus 5.5 (claude-opus-5-5, adaptive thinking, effort high), Message Batches API, batch ega25b (msgbatch_01A7n5YuXubUPNK1zsepi8RG), 2026-11-13; not yet checked by eye.
\scanpage{25/167}
\mspage{\hand{2}}

\[ v^1 : E\otimes\underline{O}_{V^1} \longrightarrow \underline{\Omega}^1_{X/Y}\otimes\underline{O}_{V^1} \]
\xout{\unsure{en termes duquel}} (déduit  lui-même de l'homomorphisme de modules
localement libres sur $V^o$
\[ v : \underline{E\otimes O}_{V^o} \longrightarrow \underline{\Omega}^1_{X/Y}\otimes \underline{O}_{V^o} \quad , \]
en termes duquel  $V^1$ a été défini comme\inserted{le} sous-préschéma de $V^o$ sur lequel
le rang de v est strictement \hand{\unsure{$\leq$}} m-1 = rang $E\otimes\underline{O}_{V^o}$ - 1). Pour la définition
du premier isomorphisme, cf page 5, pour celle du second, df compléments
différentiels (ci joints) à `` variation du rang d'un homomorphisme de
Modules loc libres''. On obtient donc une suite exacte:
\[ 0 \to \underline{E}\otimes\underline{O}_{V^1} \xrightarrow{\ i\ } \underline{N}_{V^1/X} \xrightarrow{\ \ j\ \ } \underline{\mathrm{Hom}}_{\underline{O}_{V^1}}(W,K) \longrightarrow 0 \qquad . \]
Il faut alors vérifier \xout{\unsure{1\textsuperscript{o}) que}}
1\textsuperscript{o}) que
\[ \hand{\mu}i : \underline{E}\otimes\underline{O}_{V^1} \longrightarrow \underline{\Omega}^1_{X/Y}\otimes\underline{O}_{V^1} \]
n'est autre que $v^1$ , ce qui implique alors , grâce à l'hypothèse que
$v^1$ est strictement de rang m-1, que \xout{\unsure{V(u)}} $\hand{\mu}i$ définit par passage au
quotient
\[ \hand{/\mu'} : \quad \underline{\mathrm{Hom}}\,(W,K) \longrightarrow W = \mathrm{Coker}\ v^1 \ , \]
et que V(u) est aussi le sous-préschéma de $V^1$ qui rend $\hand{\mu'}$ strictement
de rang \hand{\unsure{$\leq$}}rang W \deleted{\illegible{1}} ; d'autre part, il' \xout{peut aussi s'interpréter} faut véri-%
fier

2\textsuperscript{o}) que $\hand{\mu'}$ n'est autre que l'homomorphisme déduit de la forme
quadratique \xout{\illegible{1}} M sur W à valeurs dans K, du bas de page 15. Il en résulte
aussitôt que V($\hand{\mu'}$) = V(\over{\illegible{1}}{D}) \xout{x}, où D est le discriminant de M, ce \deleted{\unsure{qui}} qui
achève la démonstration.

La vérification de 1\textsuperscript{o}) et 2\textsuperscript{o}) (où peut-être il y a à introduire un
signe ?) ne peut être que essentiellement méchanique, et est laissée à
Dieudonné.

\margin{167}
